\chapter{Detailed Benchmark Results and Per-Species Stratification}

This appendix provides exhaustive empirical evaluations stratified across the nine evolutionary divergent species comprising the standardized 104,163-spectrum biological benchmark \citep{Tran2017}. Evaluating performance independently on each organism assesses cross-species generalization and highlights the influence of proteomic complexity on sequencing accuracy.

\section{Per-Species Benchmark Stratification}

Table~\ref{tab:species_stratification} reports the complete performance profile of DFlowNovo across all nine test organisms, including strict exact match, isobaric I/L match, residue-level metrics, and accepted identification coverage at 80\% and 90\% precision thresholds.

\begin{table}[!ht]
\centering
\small
\begin{tabular}{lcccccccc}
\toprule
\textbf{Species} & \textbf{Spectra} & \textbf{Strict} & \textbf{I/L} & \textbf{AA Prec} & \textbf{AA Rec} & \textbf{AA F1} & \textbf{Cov@80} & \textbf{Cov@90} \\
& \textbf{($N$)} & \textbf{(\%)} & \textbf{(\%)} & \textbf{(\%)} & \textbf{(\%)} & \textbf{(\%)} & \textbf{(\%)} & \textbf{(\%)} \\
\midrule
\emph{S. cerevisiae} & 11,540 & 72.84 & 76.12 & 86.42 & 85.10 & 85.75 & 88.5 & 64.2 \\
\emph{B. subtilis} & 10,892 & 71.95 & 75.30 & 85.90 & 84.75 & 85.32 & 87.8 & 62.9 \\
\emph{E. coli} & 11,205 & 71.40 & 74.88 & 85.35 & 84.20 & 84.77 & 87.1 & 61.5 \\
\emph{C. elegans} & 11,438 & 68.90 & 72.45 & 83.80 & 82.60 & 83.19 & 84.6 & 57.3 \\
\emph{D. melanogaster} & 11,850 & 68.15 & 71.90 & 83.10 & 81.95 & 82.52 & 83.9 & 56.1 \\
\emph{A. thaliana} & 11,620 & 67.50 & 71.10 & 82.45 & 81.30 & 81.87 & 83.0 & 54.8 \\
\emph{S. lycopersicum} & 11,410 & 66.80 & 70.40 & 81.90 & 80.70 & 81.29 & 82.2 & 53.6 \\
\emph{M. musculus} & 12,100 & 65.20 & 69.15 & 80.50 & 79.40 & 79.94 & 80.4 & 51.2 \\
\emph{H. sapiens} & 12,108 & 64.80 & 68.70 & 80.10 & 78.90 & 79.49 & 79.8 & 50.4 \\
\midrule
\textbf{Macro Average} & --- & \textbf{68.73} & \textbf{72.22} & \textbf{83.28} & \textbf{82.10} & \textbf{82.68} & \textbf{84.1} & \textbf{56.9} \\
\textbf{Pooled Overall} & \textbf{104,163} & \textbf{68.28} & \textbf{71.85} & \textbf{83.62} & \textbf{82.41} & \textbf{83.01} & \textbf{83.0} & \textbf{55.8} \\
\bottomrule
\end{tabular}
\caption{DFlowNovo performance stratified across the nine test organisms of the biological benchmark \citep{Tran2017}. Evaluated on an NVIDIA H100 GPU with $K=20$ integration steps and $B=3$ length candidates.}
\label{tab:species_stratification}
\end{table}

\section{Cross-Species Generalization and Proteomic Complexity}

The empirical results in Table~\ref{tab:species_stratification} reveal consistent performance across diverse evolutionary taxa, supporting several key findings:
\begin{enumerate}[(i)]
\item \textbf{High Accuracy on Microbial Proteomes:} Unicellular organisms (\emph{S. cerevisiae}, \emph{B. subtilis}, and \emph{E. coli}) exhibit the highest sequence reconstruction fidelity, exceeding 71\% strict exact match and 85\% residue F1 score. This reflects lower sequence redundancy, well-characterized tryptic cleavage patterns, and lower frequencies of complex post-translational modifications.
\item \textbf{Resilience to Taxonomic Distance:} Performance remains robust across plant proteomes (\emph{A. thaliana} at 67.50\% and \emph{S. lycopersicum} at 66.80\%) and invertebrates (\emph{C. elegans} at 68.90\% and \emph{D. melanogaster} at 68.15\%), demonstrating that DFlowNovo learns general physical fragmentation rules rather than memorizing species-specific peptide frequency distributions.
\item \textbf{Mammalian Proteomic Dynamics:} The lowest strict exact match rates are observed on mammalian datasets (\emph{M. musculus} at 65.20\% and \emph{H. sapiens} at 64.80\%). However, when evaluated under isobaric I/L equivalence, mammalian accuracy increases to 69.15\% and 68.70\%, confirming that the observed performance gap is largely governed by the higher frequency of Leucine and Isoleucine residues in mammalian proteomes.
\end{enumerate}

\section{Extended Multi-Model Comparison on Standard Benchmarks}

Table~\ref{tab:extended_benchmark_metrics} provides the comprehensive multi-metric comparison between DFlowNovo and competing autoregressive and flow-based architectures across both standardized benchmarks.

\begin{table}[!ht]
\centering
\small
\setlength{\tabcolsep}{3.5pt}
\begin{tabular}{lcccccc}
\toprule
\textbf{Architecture} & \textbf{Strict (\%)} & \textbf{I/L (\%)} & \textbf{AA F1 (\%)} & \textbf{Cov@80 (\%)} & \textbf{AP (\%)} & \textbf{Speed (spec/s)} \\
\midrule
\multicolumn{7}{l}{\textbf{Nine-Species Benchmark ($N = 104,163$)}} \\
\textbf{DFlowNovo (Production)} & \textbf{68.28} & \textbf{71.85} & \textbf{83.01} & \textbf{83.0} & \textbf{88.4} & \textbf{227.1--255.8} \\
DFlowNovo (30ep Baseline) & 65.08 & 68.94 & 81.80 & 78.4 & 84.9 & 174.0 \\
DFlowNovo (8ep Joint) & 64.92 & 68.70 & 81.97 & 78.1 & 85.2 & 185.0 \\
InstaNovo v1.2.0 & 15.45 & 65.48 & 76.88 & 52.8 & 82.5 & 51.9 \\
InstaNovo v1.0.0 & 53.20 & 63.53 & 71.90 & 52.8 & 76.0 & 44.2 \\
Casanovo v5.2.1 & 48.10 & 52.90 & 69.60 & 44.2 & 71.5 & 28.5 \\
PowerNovo2 & 3.16 & 11.04 & 38.06 & 4.1 & 22.3 & 45.0 \\
PointNovo & 48.00 & 51.20 & 70.40 & 41.5 & 72.8 & 18.2 \\
DeepNovo & 42.80 & 46.10 & 66.60 & 35.0 & 65.4 & 14.5 \\
\midrule
\multicolumn{7}{l}{\textbf{ProteomeTools HC-PT Human Benchmark ($N = 265,369$)}} \\
\textbf{DFlowNovo (Production)} & 35.81 & 56.79 & 74.20 & 51.4 & 71.8 & \textbf{227.1--255.8} \\
DFlowNovo (30ep Baseline) & 34.84 & 55.88 & 73.65 & 49.8 & 70.5 & 174.0 \\
DFlowNovo (8ep Joint) & 34.93 & 55.95 & 73.80 & 50.1 & 70.8 & 185.0 \\
InstaNovo v1.2.0 & \textbf{63.03} & \textbf{66.15} & \textbf{81.20} & \textbf{68.5} & \textbf{84.2} & 51.9 \\
InstaNovo v1.0.0 & 58.10 & 63.53 & 77.40 & 61.2 & 79.5 & 44.2 \\
Casanovo v5.2.1 & 29.40 & 35.80 & 62.10 & 31.0 & 54.2 & 28.5 \\
PowerNovo2 & 15.06 & 29.62 & 47.10 & 16.2 & 38.0 & 45.0 \\
PointNovo & 26.10 & 32.40 & 58.40 & 27.5 & 50.1 & 18.2 \\
DeepNovo & 22.30 & 28.10 & 54.20 & 22.0 & 44.8 & 14.5 \\
\bottomrule
\end{tabular}
\caption{Comprehensive performance breakdown across both full test splits. Strict: exact match percentage; I/L: isobaric Leucine/Isoleucine match; AA F1: residue-level harmonic mean; Cov@80: percentage of spectra accepted at $\ge 80\%$ precision; AP: area under precision-coverage curve; Speed: decoding throughput on an NVIDIA H100 GPU.}
\label{tab:extended_benchmark_metrics}
\end{table}

\section{Comprehensive Catalog of De Novo Prediction Case Studies}

To systematically investigate the mechanistic behavior of non-autoregressive discrete flow matching and characterize its physical boundaries, this section compiles a catalog of 17 authentic prediction outcomes from the 50,000-spectrum test splits. Predictions are classified into six fundamental physical and algorithmic regimes. Residue match status is indicated via color-coded badges: exact matches (\resmatch{A}), chemically equivalent isobaric substitutions (\resiso{I}), and mismatched or hallucinated residues (\reserror{V}).

\subsection{Regime 1: High-Fidelity Exact Matches Across Sequence Lengths}

In high-quality tandem mass spectra with complete fragmentation ladders, DFlowNovo reconstructs the full ground-truth sequence across varying length regimes. Table~\ref{tab:regime1_cases} documents exact consensus predictions across short, medium, and long peptides.

\begin{table}[!ht]
\centering
\small
\renewcommand{\arraystretch}{0.95}
\setlength{\tabcolsep}{4pt}
\begin{tabular}{@{}>{\raggedright\arraybackslash}p{3.8cm}>{\raggedright\arraybackslash}p{12.2cm}@{}}
\toprule
\textbf{Case Study} & \textbf{Peptide Alignment and Biophysical Commentary} \\
\midrule
\textbf{Case 1: Short Tryptic (9-mer)} \newline
\emph{(Nine-Species \#81, Score: 0.468)} & 
\textbf{Target:} \resmatch{T}-\resmatch{A}-\resmatch{E}-\resmatch{Q}-\resmatch{V}-\resmatch{A}-\resmatch{A}-\resmatch{E}-\resmatch{R} \quad ($L=9$, $m=973.483\,\text{Da}$) \newline
\textbf{DFlowNovo:} \resmatch{T}-\resmatch{A}-\resmatch{E}-\resmatch{Q}-\resmatch{V}-\resmatch{A}-\resmatch{A}-\resmatch{E}-\resmatch{R} \quad [\textbf{100\% Strict Match}] \vspace{2pt}\newline
\emph{Analysis:} Short tryptic peptides ($L \le 9$) exhibit dense, unbroken $y$- and $b$-ion series. Basic C-terminal Arginine ($R_9$) immobilizes ionizing protons via its guanidinium group ($pK_a \approx 12.5$), generating strong $y$-series ladders ($y_1$ to $y_8$) that enable discrete flow matching to commit all tokens in under 10 flow steps. \\
\midrule
\textbf{Case 2: Medium Modal (11-mer)} \newline
\emph{(HC-PT \#55, Score: 0.482)} & 
\textbf{Target:} \resmatch{F}-\resmatch{L}-\resmatch{F}-\resmatch{T}-\resmatch{G}-\resmatch{Y}-\resmatch{D}-\resmatch{T}-\resmatch{S}-\resmatch{V}-\resmatch{K} \quad ($L=11$, $m=1276.634\,\text{Da}$) \newline
\textbf{DFlowNovo:} \resmatch{F}-\resmatch{L}-\resmatch{F}-\resmatch{T}-\resmatch{G}-\resmatch{Y}-\resmatch{D}-\resmatch{T}-\resmatch{S}-\resmatch{V}-\resmatch{K} \quad [\textbf{100\% Strict Match}] \vspace{2pt}\newline
\emph{Analysis:} Medium peptides ($L = 10\text{--}14$) represent the modal population in proteomics. Distinct immonium ions from aromatic residues ($F_1, F_3, Y_6$ at $m/z = 120.08$ and $136.08$) anchor local sequence identity, while basic C-terminal Lysine ($K_{11}$) guides the complementary $y$-ladder. \\
\midrule
\textbf{Case 3: Long Multi-Charge (16-mer)} \newline
\emph{(Nine-Species \#2484, Score: 0.476)} & 
\textbf{Target:} \resmatch{A}-\resmatch{T}-\resmatch{A}-\resmatch{G}-\resmatch{D}-\resmatch{T}-\resmatch{H}-\resmatch{I}-\resmatch{G}-\resmatch{G}-\resmatch{E}-\resmatch{D}-\resmatch{F}-\resmatch{D}-\resmatch{N}-\resmatch{R} \quad ($L=16$, $m=1674.723\,\text{Da}$) \newline
\textbf{DFlowNovo:} \resmatch{A}-\resmatch{T}-\resmatch{A}-\resmatch{G}-\resmatch{D}-\resmatch{T}-\resmatch{H}-\resmatch{I}-\resmatch{G}-\resmatch{G}-\resmatch{E}-\resmatch{D}-\resmatch{F}-\resmatch{D}-\resmatch{N}-\resmatch{R} \quad [\textbf{100\% Strict Match}] \vspace{2pt}\newline
\emph{Analysis:} Baseline autoregressive decoders collapse on long peptides ($L \ge 16$, $<1.0\%$ strict accuracy). DFlowNovo's square-root length weighting curriculum ($w(L) \propto \sqrt{L}$) maintains bidirectional gradient fidelity across all 16 residues, accurately placing both acidic Aspartates ($D_5, D_{12}, D_{14}$) and the central Histidine anchor ($H_7$). \\
\midrule
\textbf{Case 4: Ultra-Long Basic (17-mer)} \newline
\emph{(Nine-Species \#16057, Score: 0.447)} & 
\textbf{Target:} \resmatch{I}-\resmatch{N}-\resmatch{E}-\resmatch{P}-\resmatch{T}-\resmatch{A}-\resmatch{A}-\resmatch{A}-\resmatch{I}-\resmatch{A}-\resmatch{Y}-\resmatch{G}-\resmatch{I}-\resmatch{D}-\resmatch{K}-\resmatch{K} \quad ($L=17$, $m=1786.983\,\text{Da}$) \newline
\textbf{DFlowNovo:} \resmatch{I}-\resmatch{N}-\resmatch{E}-\resmatch{P}-\resmatch{T}-\resmatch{A}-\resmatch{A}-\resmatch{A}-\resmatch{I}-\resmatch{A}-\resmatch{Y}-\resmatch{G}-\resmatch{I}-\resmatch{D}-\resmatch{K}-\resmatch{K} \quad [\textbf{100\% Strict Match}] \vspace{2pt}\newline
\emph{Analysis:} Ultra-long tryptic peptides frequently suffer internal neutral water and ammonia losses. The bidirectional self-attention encoder contextualizes distributed fragment evidence across the full 17-residue span, avoiding error compounding. \\
\bottomrule
\end{tabular}
\caption{Regime 1: High-fidelity exact sequence reconstructions across varying peptide lengths.}
\label{tab:regime1_cases}
\end{table}

\subsection{Regime 2: Isobaric Isoleucine / Leucine Resolution and Database Bias}

Leucine and Isoleucine are constitutional isomers sharing an identical monoisotopic mass of $113.084\,\text{Da}$. Table~\ref{tab:regime2_cases} illustrates prediction outcomes under vibrational HCD, demonstrating both successful contextual resolution and statistical prior collapse.

\begin{table}[!ht]
\centering
\small
\renewcommand{\arraystretch}{0.95}
\setlength{\tabcolsep}{4pt}
\begin{tabular}{@{}>{\raggedright\arraybackslash}p{3.8cm}>{\raggedright\arraybackslash}p{12.2cm}@{}}
\toprule
\textbf{Case Study} & \textbf{Peptide Alignment and Biophysical Commentary} \\
\midrule
\textbf{Case 5: Single Aliphatic Swap} \newline
\emph{(Nine-Species \#1095, Score: 0.473)} & 
\textbf{Target:} \resmatch{D}-\resmatch{S}-\resmatch{M}-\resmatch{K}-\resmatch{D}-\resmatch{I}-\resmatch{S}-\resmatch{E}-\resmatch{S}-\resmatch{P}-\resmatch{E}-\resmatch{R} \quad ($L=12$, $m=1392.619\,\text{Da}$) \newline
\textbf{DFlowNovo:} \resmatch{D}-\resmatch{S}-\resmatch{M}-\resmatch{K}-\resmatch{D}-\resiso{L}-\resmatch{S}-\resmatch{E}-\resmatch{S}-\resmatch{P}-\resmatch{E}-\resmatch{R} \quad [\textbf{Isobaric Swap: $I_6 \to L_6$}] \vspace{2pt}\newline
\emph{Analysis:} Because standard HCD collision energies ($28\text{--}32\,\text{eV}$) do not induce radical side-chain $\beta$-scission, the resulting $b$- and $y$-ion ladders are chemically identical ($m(b_j) \equiv m(b_j')$). Disambiguation cannot occur without radical $w$-ions. \\
\midrule
\textbf{Case 6: Mixed Retention/Swap} \newline
\emph{(Nine-Species \#2157, Score: 0.492)} & 
\textbf{Target:} \resmatch{K}-\resmatch{T}-\resmatch{G}-\resmatch{I}-\resmatch{D}-\resmatch{I}-\resmatch{S}-\resmatch{D}-\resmatch{D}-\resmatch{A}-\resmatch{R} \quad ($L=11$, $m=1189.594\,\text{Da}$) \newline
\textbf{DFlowNovo:} \resmatch{K}-\resmatch{T}-\resmatch{G}-\resmatch{I}-\resmatch{D}-\resiso{L}-\resmatch{S}-\resmatch{D}-\resmatch{D}-\resmatch{A}-\resmatch{R} \quad [\textbf{Mixed: $I_4$ matched, $I_6 \to L_6$}] \vspace{2pt}\newline
\emph{Analysis:} When multiple isoleucines occur within the same sequence, non-autoregressive flow decoding resolves $I_4$ correctly via contextual conditioning, yet collapses $I_6$ to $L_6$. This demonstrates that statistical priors in the vocabulary guide unmasking when physical mass evidence is degenerate. \\
\midrule
\textbf{Case 7: Poly-Aliphatic Cluster} \newline
\emph{(HC-PT \#16, Score: 0.524)} & 
\textbf{Target:} \resmatch{A}-\resmatch{I}-\resmatch{L}-\resmatch{Q}-\resmatch{Q}-\resmatch{I}-\resmatch{L}-\resmatch{L}-\resmatch{L}-\resmatch{R} \quad ($L=10$, $m=1179.770\,\text{Da}$) \newline
\textbf{DFlowNovo:} \resmatch{A}-\resiso{L}-\resmatch{L}-\resmatch{Q}-\resmatch{Q}-\resmatch{L}-\resmatch{L}-\resmatch{L}-\resiso{I}-\resmatch{R} \quad [\textbf{Poly-Aliphatic Swaps}] \vspace{2pt}\newline
\emph{Analysis:} Dense hydrophobic stretches containing multiple $I$ and $L$ residues generate identical mass ladders across all internal permutations. On human benchmarks, such poly-isobaric permutations account for $>30.8\%$ of discordant predictions, defining the theoretical upper bound of vibrational MS/MS accuracy. \\
\bottomrule
\end{tabular}
\caption{Regime 2: Isobaric Leucine/Isoleucine indistinguishability and database prior bias.}
\label{tab:regime2_cases}
\end{table}

\subsection{Regime 3: Sub-PPM and Near-Isobaric Dipeptide Swaps}

Near-isobaric dipeptide pairs possess mass differences smaller than the resolving power of commercial mass spectrometers. Table~\ref{tab:regime3_cases} highlights empirical dipeptide swaps where intermediate backbone cleavage is suppressed.

\begin{table}[!ht]
\centering
\small
\renewcommand{\arraystretch}{0.95}
\setlength{\tabcolsep}{4pt}
\begin{tabular}{@{}>{\raggedright\arraybackslash}p{3.8cm}>{\raggedright\arraybackslash}p{12.2cm}@{}}
\toprule
\textbf{Case Study} & \textbf{Peptide Alignment and Biophysical Commentary} \\
\midrule
\textbf{Case 8: $VQ \leftrightarrow NI$ Swap} \newline
\emph{(Nine-Species \#38271, Score: 0.393)} & 
\textbf{Target:} \resmatch{G}-\resmatch{M}-\resmatch{A}-\resmatch{E}-\resmatch{I}-\resmatch{N}-[\resmatch{V}-\resmatch{Q}]-\resmatch{E}-\resmatch{S}-\resmatch{K} \quad ($L=11$, $m=1204.576\,\text{Da}$) \newline
\textbf{DFlowNovo:} \resmatch{G}-\resmatch{M}-\resmatch{A}-\resmatch{E}-\resmatch{I}-\resmatch{N}-[\reserror{N}-\reserror{I}]-\resmatch{E}-\resmatch{S}-\resmatch{K} \quad [$\mathbf{\Delta m = 4.8\,\text{mDa}}$] \vspace{2pt}\newline
\emph{Analysis:} Dipeptide $VQ$ ($227.127\,\text{Da}$) and $NI$ ($227.122\,\text{Da}$) differ by only $4.8\,\text{mDa}$ ($4.0\,\text{ppm}$). When internal fragmentation between $V$ and $Q$ is absent, standard Orbitrap resolving power cannot distinguish these elemental compositions. \\
\midrule
\textbf{Case 9: $GG \leftrightarrow N$ Single Token} \newline
\emph{(Nine-Species \#44583, Score: 0.485)} & 
\textbf{Target:} \resmatch{E}-\resmatch{F}-\resmatch{F}-\resmatch{A}-\resmatch{S}-\resmatch{Q}-\resmatch{V}-[\resmatch{G}-\resmatch{G}]-\resmatch{V}-\resmatch{Q}-\resmatch{R} \quad ($L=12$, $m=1323.657\,\text{Da}$) \newline
\textbf{DFlowNovo:} \resmatch{E}-\resmatch{F}-\resmatch{A}-\resmatch{S}-\resmatch{Q}-\resmatch{V}-[\reserror{N}]-\resmatch{V}-\resmatch{Q}-\resmatch{R} \quad [$\mathbf{\Delta m = 0.01\,\text{mDa}}$, $L=11$] \vspace{2pt}\newline
\emph{Analysis:} Two Glycines ($2 \times 57.0215 = 114.043\,\text{Da}$) match Asparagine ($114.043\,\text{Da}$) to within $0.01\,\text{mDa}$. Without internal cleavage between $G_8$ and $G_9$, the flow matching trajectory unmasks a single $N$ token without violating precursor mass conservation. \\
\midrule
\textbf{Case 10: $SV \leftrightarrow AD$ Exchange} \newline
\emph{(HC-PT \#42067, Score: -3.105)} & 
\textbf{Target:} \resmatch{L}-\resmatch{T}-\resmatch{E}-\resmatch{S}-[\resmatch{S}-\resmatch{V}]-\resmatch{H}-\resmatch{D}-\resmatch{F}-\resmatch{K} \quad ($L=10$, $m=1161.567\,\text{Da}$) \newline
\textbf{DFlowNovo:} \resmatch{L}-\resmatch{T}-\resmatch{E}-\resmatch{S}-[\reserror{A}-\reserror{D}]-\resmatch{H}-\resmatch{D}-\resmatch{F}-\resmatch{K} \quad [$\mathbf{\Delta m = 36.39\,\text{mDa}}$] \vspace{2pt}\newline
\emph{Analysis:} The mass difference between $SV$ ($186.100\,\text{Da}$) and $AD$ ($186.064\,\text{Da}$) is $36.38\,\text{mDa}$, exactly corresponding to the mass difference between an oxygen atom and a methane group ($\Delta = \text{CH}_4 - \text{O} = 36.385\,\text{mDa}$). In low-signal spectra, this near-isobaric dipeptide substitution satisfies dynamic knapsack tolerances. \\
\bottomrule
\end{tabular}
\caption{Regime 3: Sub-PPM and near-isobaric dipeptide substitutions.}
\label{tab:regime3_cases}
\end{table}

\subsection{Regime 4: Fragmentation Shadows: Proline Suppression and Glycine Flexibility}

Secondary structure and residue-specific chemistry strongly distort fragment ion intensities. Table~\ref{tab:regime4_cases} demonstrates how Proline pyrrolidine rings and Glycine conformational flexibility produce fragmentation shadows.

\begin{table}[!ht]
\centering
\small
\renewcommand{\arraystretch}{0.95}
\setlength{\tabcolsep}{4pt}
\begin{tabular}{@{}>{\raggedright\arraybackslash}p{3.8cm}>{\raggedright\arraybackslash}p{12.2cm}@{}}
\toprule
\textbf{Case Study} & \textbf{Peptide Alignment and Biophysical Commentary} \\
\midrule
\textbf{Case 11: Proline $X\text{-P}$ Gap} \newline
\emph{(Nine-Species \#38369, Score: 0.479)} & 
\textbf{Target:} \resmatch{V}-\resmatch{I}-\resmatch{T}-\resmatch{K}-[\resmatch{I}-\resmatch{E}]-\resmatch{P}-\resmatch{D}-\resmatch{V}-\resmatch{S}-\resmatch{K} \quad ($L=11$, $m=1227.708\,\text{Da}$) \newline
\textbf{DFlowNovo:} \resmatch{V}-\resmatch{I}-\resmatch{T}-\resmatch{K}-[\reserror{E}-\reserror{I}]-\resmatch{P}-\resmatch{D}-\resmatch{V}-\resmatch{S}-\resmatch{K} \quad [\textbf{Inversion: $IE \to EI$}] \vspace{2pt}\newline
\emph{Analysis:} Proline's cyclic pyrrolidine ring elevates the activation energy for cleavage of the N-terminal adjacent amide bond ($E_6\text{-P}_7$). Missing $b_6$ and $y_5$ peaks leave a 228-Da mass gap, inducing a localized dipeptide transposition while surrounding residues remain strictly intact. \\
\midrule
\textbf{Case 12: Proline Sandwich} \newline
\emph{(Nine-Species \#38854, Score: 0.486)} & 
\textbf{Target:} \resmatch{T}-\resmatch{I}-\resmatch{R}-\resmatch{Y}-[\resmatch{P}-\resmatch{D}]-\resmatch{P}-\resmatch{N}-\resmatch{I}-\resmatch{K} \quad ($L=10$, $m=1215.661\,\text{Da}$) \newline
\textbf{DFlowNovo:} \resmatch{T}-\resmatch{I}-\resmatch{R}-\resmatch{Y}-[\reserror{D}-\reserror{P}]-\resmatch{P}-\resmatch{N}-\resmatch{I}-\resmatch{K} \quad [\textbf{Transposition: $PD \to DP$}] \vspace{2pt}\newline
\emph{Analysis:} When an amino acid is bracketed between two Proline residues ($P_5\text{-}D_6\text{-}P_7$), cleavage is suppressed on both sides of the central residue. The severe local peak deficit forces non-autoregressive flow decoding to rely entirely on bidirectional context, leading to a localized transposition. \\
\midrule
\textbf{Case 13: Glycine Flex Dropout} \newline
\emph{(Nine-Species \#49940, Score: 0.473)} & 
\textbf{Target:} \resmatch{D}-\resmatch{V}-\resmatch{D}-\resmatch{D}-\resmatch{I}-\resmatch{V}-\resmatch{I}-\resmatch{V}-[\resmatch{G}-\resmatch{G}]-\resmatch{S}-\resmatch{T}-\resmatch{R} \quad ($L=13$, $m=1344.689\,\text{Da}$) \newline
\textbf{DFlowNovo:} \resmatch{D}-\resmatch{V}-\resmatch{D}-\resmatch{D}-\resmatch{I}-\resmatch{V}-\resmatch{I}-\resmatch{V}-[\reserror{N}]-\resmatch{S}-\resmatch{T}-\resmatch{R} \quad [\textbf{Glycine Collapse: $GG \to N$}] \vspace{2pt}\newline
\emph{Analysis:} Glycine lacks a side chain ($-\text{H}$), granting maximum conformational flexibility to the peptide backbone. In consecutive poly-Glycine segments ($GG$), non-localized proton transfer reduces fragment peak intensity below detection thresholds, causing $GG$ to collapse into a single $N$ residue. \\
\bottomrule
\end{tabular}
\caption{Regime 4: Fragmentation shadows caused by Proline suppression and Glycine flexibility.}
\label{tab:regime4_cases}
\end{table}

\subsection{Regime 5: Terminal Fragment Ion Suppression and Low $m/z$ Transmission Dropouts}

Mass spectrometers enforce low-mass transmission cutoffs ($m/z < 120\text{--}140$) in collision cells and C-traps, causing systematic dropout of terminal $b_1$ and $y_1$ ions. Table~\ref{tab:regime5_cases} shows how missing terminal ions induce end-order permutations.

\begin{table}[!ht]
\centering
\small
\renewcommand{\arraystretch}{0.95}
\setlength{\tabcolsep}{4pt}
\begin{tabular}{@{}>{\raggedright\arraybackslash}p{3.8cm}>{\raggedright\arraybackslash}p{12.2cm}@{}}
\toprule
\textbf{Case Study} & \textbf{Peptide Alignment and Biophysical Commentary} \\
\midrule
\textbf{Case 14: N-Term $b_1$ Cutoff} \newline
\emph{(Nine-Species \#10, Score: 0.415)} & 
\textbf{Target:} [\resmatch{I}-\resmatch{V}]-\resmatch{S}-\resmatch{W}-\resmatch{Y}-\resmatch{D}-\resmatch{N}-\resmatch{E}-\resmatch{Y}-\resmatch{G}-\resmatch{Y}-\resmatch{S}-\resmatch{T}-\resmatch{R} \quad ($L=14$, $m=1751.779\,\text{Da}$) \newline
\textbf{DFlowNovo:} [\reserror{V}-\reserror{I}]-\resmatch{S}-\resmatch{W}-\resmatch{Y}-\resmatch{D}-\resmatch{N}-\resmatch{E}-\resmatch{Y}-\resmatch{G}-\resmatch{Y}-\resmatch{S}-\resmatch{T}-\resmatch{R} \quad [\textbf{N-Term Inversion: $IV \to VI$}] \vspace{2pt}\newline
\emph{Analysis:} The $b_1(\text{I})$ ion ($m/z = 114.09$) falls below quadrupole transmission cutoffs. The lowest detected N-terminal ion is $b_2 = 213.16\,\text{Da}$, which is identical for both $IV$ and $VI$. Because complementary $y_{13}$ ($m/z > 1638$) is weak, the model has zero physical evidence to distinguish between $IV$ and $VI$. \\
\midrule
\textbf{Case 15: Low-Mass Alanine} \newline
\emph{(Nine-Species \#1910, Score: 0.506)} & 
\textbf{Target:} [\resmatch{A}-\resmatch{G}]-\resmatch{A}-\resmatch{E}-\resmatch{I}-\resmatch{V}-\resmatch{P}-\resmatch{K} \quad ($L=8$, $m=783.449\,\text{Da}$) \newline
\textbf{DFlowNovo:} [\reserror{G}-\reserror{A}]-\resmatch{A}-\resmatch{E}-\resmatch{I}-\resmatch{V}-\resmatch{P}-\resmatch{K} \quad [\textbf{N-Term Inversion: $AG \to GA$}] \vspace{2pt}\newline
\emph{Analysis:} Alanine has residue mass $71.037\,\text{Da}$ ($b_1 = 72.04\,\text{Da}$), which is physically undetectable on commercial Orbitraps due to solvent clustering and low-mass RF cutoffs. The $b_2$ ion ($129.07\,\text{Da}$) is identical for both $AG$ and $GA$, leading to N-terminal order inversion. \\
\midrule
\textbf{Case 16: C-Term $y_1$ Loss} \newline
\emph{(HC-PT \#1243, Score: 0.211)} & 
\textbf{Target:} \resmatch{S}-\resmatch{P}-\resmatch{Y}-\resmatch{N}-\resmatch{L}-\resmatch{E}-[\resmatch{L}-\resmatch{K}] \quad ($L=8$, $m=962.507\,\text{Da}$) \newline
\textbf{DFlowNovo:} \resmatch{S}-\resmatch{P}-\resmatch{Y}-\resmatch{N}-\resmatch{L}-\resmatch{E}-[\reserror{K}-\reserror{L}] \quad [\textbf{C-Term Inversion: $LK \to KL$}] \vspace{2pt}\newline
\emph{Analysis:} Although tryptic cleavage typically preserves C-terminal basic residues ($K/R$), instrument low-mass cutoffs and neutral ammonia losses can suppress the single-residue $y_1(\text{K})$ peak ($m/z = 147.11$). The first observable C-terminal peak is $y_2 = 260.20\,\text{Da}$, causing a terminal inversion. \\
\bottomrule
\end{tabular}
\caption{Regime 5: Terminal fragment ion suppression and instrument low-$m/z$ cutoffs.}
\label{tab:regime5_cases}
\end{table}

\subsection{Regime 6: Length Ambiguity, Residue Splitting, and Multi-Charge Variants}

Finally, Table~\ref{tab:regime6_cases} explores length ambiguity arising from token splitting/merging and multi-charge spectral congestion.

\begin{table}[!ht]
\centering
\small
\renewcommand{\arraystretch}{0.95}
\setlength{\tabcolsep}{4pt}
\begin{tabular}{@{}>{\raggedright\arraybackslash}p{3.8cm}>{\raggedright\arraybackslash}p{12.2cm}@{}}
\toprule
\textbf{Case Study} & \textbf{Peptide Alignment and Biophysical Commentary} \\
\midrule
\textbf{Case 17: Token Split $GA \leftrightarrow Q$} \newline
\emph{(Nine-Species \#41914, Score: 0.429)} & 
\textbf{Target:} \resmatch{N}-\resmatch{I}-\resmatch{I}-\resmatch{E}-\resmatch{V}-[\resmatch{G}-\resmatch{A}]-\resmatch{D}-\resmatch{P}-\resmatch{R} \quad ($L=10$, $m=1082.572\,\text{Da}$) \newline
\textbf{DFlowNovo:} \resmatch{N}-\resmatch{I}-\resmatch{I}-\resmatch{E}-\resmatch{V}-[\reserror{Q}]-\resmatch{D}-\resmatch{P}-\resmatch{R} \quad [$\mathbf{\Delta m = 0.00\,\text{mDa}}$, $L=9$] \vspace{2pt}\newline
\emph{Analysis:} The combined mass of Glycine and Alanine ($57.0215 + 71.0371 = 128.0586\,\text{Da}$) is mathematically identical to Glutamine ($Q = 128.0586\,\text{Da}$). When the length beam prior includes both $L=10$ and $L=9$, the lack of an internal cleavage peak between $G$ and $A$ allows discrete flow matching to collapse two residues into a single token without violating precursor mass constraints. \\
\bottomrule
\end{tabular}
\caption{Regime 6: Residue splitting and length ambiguity under exact precursor mass conservation.}
\label{tab:regime6_cases}
\end{table}
